\chapter{Conclusão}
\label{cap:conclusao}

Esta dissertação investigou se criptogramas produzidos por Ascon-AEAD128 e GIFT-COFB podem ser distinguidos por classificadores de aprendizado de máquina em cenário ciphertext-only, no qual o adversário observa apenas o texto cifrado, sem acesso à chave, ao nonce, ao texto em claro ou à tag de autenticação. Para isso, foram gerados 60.000 criptogramas de 64~KB a partir de 30.000 mensagens do Standardized Project Gutenberg Corpus, cifradas de forma pareada pelos dois algoritmos com a mesma chave, o mesmo nonce e o mesmo texto em claro, e o protocolo experimental foi desenhado para bloquear os mecanismos de vazamento identificados na literatura: particionamento por chave, com 240 chaves para treinamento e validação e 60 chaves reservadas ao teste, e seleção de atributos recalculada dentro de cada dobra da validação cruzada.

Quatro caminhos experimentais foram avaliados sobre esse protocolo: descritores estatísticos clássicos (Caminho A), representação aprendida por CNN-1D (Caminho B), mapa de co-ocorrência de bigramas com CNN-2D (Caminho C) e fusão híbrida de 947 dimensões (Caminho D). Excluídos os modos de diagnóstico, o F1-macro no conjunto de teste ficou entre 0,4952 e 0,5040 em todas as combinações de representação e classificador, com os intervalos de confiança de 95\% cobrindo o valor de 0,50 esperado ao acaso, e o teste de McNemar não indicou diferença significativa entre classificadores após a correção de Bonferroni. Dois indicadores de segunda ordem apontaram na mesma direção: a busca em grade do SVM RBF não convergiu para nenhuma região do espaço de hiperparâmetros ao longo das dobras, e a estabilidade da seleção de atributos, medida pelo índice de Jaccard, caiu de 0,32 no Caminho A para 0,0013 no espaço híbrido do Caminho D.

O controle positivo delimita a leitura desses números. Submetida ao mesmo pipeline de extração, seleção e treinamento, a cifra de Vigenère com chave curta foi separada de bytes aleatórios com recall e precisão de 1,0 nos quatro classificadores. O instrumento detecta padrão quando ele existe; no par Ascon-AEAD128/GIFT-COFB, não encontrou. A resposta à pergunta de pesquisa formulada na Introdução é, portanto, negativa: sob o protocolo avaliado, não foi possível identificar, a partir exclusivamente de criptogramas, qual algoritmo de criptografia leve os gerou. Esse resultado constitui evidência empírica de indistinguibilidade prática entre os dois esquemas no cenário mais restrito da taxonomia de ameaças, e não uma prova formal de segurança, cujas garantias permanecem ancoradas nas noções de indistinguibilidade computacional apresentadas no Capítulo 2.

\section{Contribuições}

Este trabalho contribui em três frentes. A primeira é metodológica: um protocolo experimental que neutraliza, por construção, os dois mecanismos de inflação de desempenho documentados no Capítulo 3, o reaproveitamento de chave entre treinamento e teste e a seleção de características fora do laço de validação cruzada. A segunda é a evidência empírica em si: quatro representações independentes do criptograma, cobrindo descritores manuais, representação sequencial aprendida, representação de segunda ordem e a fusão das três, convergem para o desempenho de acaso, com a sensibilidade do pipeline verificada por controle positivo. A terceira é o artefato: um dataset reprodutível de 60.000 criptogramas com manifesto de geração, seeds fixas, validação por 1.089 vetores de teste oficiais por algoritmo e verificações de integridade documentadas, utilizável como benchmark por trabalhos futuros de identificação algorítmica.

\section{Limitações e Trabalhos Futuros}

Os resultados estão condicionados às escolhas do protocolo, discutidas em detalhe na Seção~\ref{sec:res-sintese}: um único comprimento de mensagem (64~KB), texto em claro proveniente de um único corpus, duas cifras de um mesmo processo de padronização e as implementações de referência em C. Trabalhos futuros podem estender a avaliação a mensagens curtas, nas quais o peso relativo da inicialização e da tag de autenticação é maior; a outros pares de esquemas AEAD leves, como Xoodyak e Romulus; a implementações otimizadas ou em hardware, nas quais artefatos de implementação podem introduzir sinal ausente das implementações de referência; e a outros corpora de texto em claro. Permanece em aberto, ainda, a execução dos controles negativos previstos no protocolo, o embaralhamento de bytes de criptogramas reais e a classe de bytes de gerador pseudoaleatório puro, que verificariam a especificidade do pipeline na ausência de qualquer estrutura.
